% ═══════════════════════════════════════════════════════════════════════
\section{Livrable 1 : une activité d'enseignement}
% ═══════════════════════════════════════════════════════════════════════

\begin{frame}{La consigne}
  \begin{Exercice}{Présenter une activité mathématique, en 2 à 3 pages}
    \small
    \begin{itemize}
      \item Classe \texttt{article}, compilée sans erreur sur Overleaf.
      \item Contenu : titre ; niveau visé, durée, intention, préalables,
        matériel ; mise en situation ; déroulement ; tâche de l'élève ou de l'étudiant ;
        corrigé et pistes d'intervention.
      \item Au moins une figure (TikZ ou image) ou un tableau.
      \item Des mathématiques en mode mathématique : formules propres,
        virgule décimale, unités avec \cmd{qty}.
      \item Un en-tête de page avec le niveau et le titre de l'activité.
    \end{itemize}
  \end{Exercice}
  \small Point de départ possible : le fichier \texttt{02-activite.tex},
  ou un squelette produit par l'IA avec le prompt suivant.
\end{frame}

\begin{frame}{Prompt type : la fiche d'activité}
  \begin{Prompt}{Fiche de présentation d'une activité}
    \footnotesize Tu es conseiller pédagogique en mathématiques
    \champ{au secondaire ou au collégial} au Québec. Rédige en LaTeX (classe article, pdfLaTeX,
    Overleaf) la présentation d'une activité de \champ{75} minutes pour
    des \champ{élèves de 3\textsuperscript{e} secondaire, ou étudiants d'un
    cours collégial} sur \champ{la relation de
    Pythagore, dans le contexte de la taille des écrans de télévision}.

    Structure : titre ; un tableau de présentation (niveau, durée,
    intention, préalables, matériel) ; mise en situation ; déroulement
    minuté ; tâche de l'élève en \champ{4} questions graduées ; corrigé
    détaillé ; pistes d'intervention pour deux erreurs fréquentes.

    Contraintes : 2 à 3 pages ; français ; une figure TikZ simple ; un
    en-tête avec fancyhdr ; siunitx avec virgule décimale pour tous les
    nombres décimaux et les unités ; extensions standard seulement.
    Donne un seul fichier .tex complet, commenté.
  \end{Prompt}
\end{frame}

\begin{frame}{Un exemple de rendu : \texttt{02-activite.tex}}
  \begin{columns}[c]
    \begin{column}{0.33\textwidth}
      \centering
      \fbox{\includegraphics[height=0.66\textheight, page=1]{02-activite.pdf}}
    \end{column}
    \begin{column}{0.33\textwidth}
      \centering
      \fbox{\includegraphics[height=0.66\textheight, page=2]{02-activite.pdf}}
    \end{column}
    \begin{column}{0.3\textwidth}
      \small
      \begin{itemize}
        \item Titre, en-tête et pied de page.
        \item Tableau de présentation (liste \texttt{description}).
        \item Figure TikZ du téléviseur.
        \item Tâche et corrigé numérotés a), b), c), d).
      \end{itemize}
    \end{column}
  \end{columns}
\end{frame}

\begin{frame}[fragile]{Deux blocs de la fiche à savoir modifier}
  \begin{columns}[T]
    \begin{column}{0.5\textwidth}
\begin{CodeTeX}[title={L'en-tête de page (préambule)}]
\usepackage{fancyhdr}
\pagestyle{fancy}
\fancyhf{}                % on efface tout
\lhead{Mathématique, 3\textsuperscript{e} secondaire}
\rhead{Activité : Pythagore au salon}
\cfoot{\thepage}          % numéro au centre
\end{CodeTeX}
    \end{column}
    \begin{column}{0.47\textwidth}
\begin{CodeTeX}[title={Une question et sa réponse (corps)}]
\begin{enumerate}[label=\alph*)]
  \item Calcule la largeur, en
    centimètres, d'un écran de
    \qty{55}{po}.
\end{enumerate}
...
$k = \dfrac{55}{\sqrt{337}}
   \approx \num{2,996}$
\end{CodeTeX}
    \end{column}
  \end{columns}
  \vspace{4pt}
  \small \texttt{[label=\textbackslash alph*)]} (extension
  \texttt{enumitem}) numérote a), b), c). Changer le niveau, le titre ou
  une valeur numérique : une seule ligne à modifier, à la main.
\end{frame}
